\chapter{Screening with a language model}
\label{ch:llm-screening}

% Author's note (verification and reproduction): session screenlit-q4
% (and screenlit-q4-top); queries q4a-q4e, q4x-q4z; check by meaning
% checks/q4.json; sample tier2-sample-q4.tsv (seed 20261127); see the
% block at the top of search.tex.  Sources: doc/bibliography.bib, keys
% Xie2026 Tran2024 Sandner2025 Krag2024 Khraisha2024 Sciurti2025
% Huotala2026 Figalova2026 Teo2023 Mateen2013.

Scholar's classifier screens records with a language model, and an agent
can screen them too (\cref{ch:procedure}).
A model's verdict is cheap, and a cheaper first pass still (titles only,
a small model, a model trusted where it is confident) would be cheaper.
This chapter asks what such a pass costs in missed studies:
\emph{does a first pass with a language model miss few relevant studies
compared with human screening?}
The claim is tested as one a procedure could lean on without a human
check: few missed studies, across reviews.

\section{Method}
\label{sec:llm-method}

The search is the one described in \cref{ch:lit-search}.
The fourth question's queries are listed below, and
\cref{tab:llm-matrix} gives their hits per database.
The support queries are q4a to q4e, the counter queries q4x to q4z; no
query was added after the check by meaning.
\begin{description}[style=nextline,font=\normalfont\bfseries]
\item[q4a, support query] {\footnotesize\ttfamily\raggedright (\textquotedbl{}large language model\textquotedbl{} OR \textquotedbl{}large language models\textquotedbl{} OR GPT OR ChatGPT) AND (\textquotedbl{}title and abstract\textquotedbl{} OR \textquotedbl{}title-and-abstract\textquotedbl{} OR \textquotedbl{}abstract screening\textquotedbl{} OR \textquotedbl{}citation screening\textquotedbl{}) AND (\textquotedbl{}systematic review\textquotedbl{} OR \textquotedbl{}literature review\textquotedbl{})\par}
\item[q4b, support query] {\footnotesize\ttfamily\raggedright (\textquotedbl{}title only\textquotedbl{} OR \textquotedbl{}titles only\textquotedbl{} OR \textquotedbl{}title-only\textquotedbl{} OR \textquotedbl{}title screening\textquotedbl{}) AND (screening) AND (abstract) AND (\textquotedbl{}systematic review\textquotedbl{})\par}
\item[q4c, support query] {\footnotesize\ttfamily\raggedright (\textquotedbl{}large language model\textquotedbl{} OR GPT OR LLM) AND (screening) AND (cascade OR \textquotedbl{}two-stage\textquotedbl{} OR uncertain OR \textquotedbl{}smaller model\textquotedbl{} OR cost OR \textquotedbl{}cost-effective\textquotedbl{}) AND (\textquotedbl{}systematic review\textquotedbl{})\par}
\item[q4d, support query] {\footnotesize\ttfamily\raggedright (\textquotedbl{}large language model\textquotedbl{} OR GPT OR LLM OR ChatGPT) AND (screening) AND (sensitivity OR recall OR missed OR \textquotedbl{}false negative\textquotedbl{} OR error) AND (human OR reviewers)\par}
\item[q4e, support query] {\footnotesize\ttfamily\raggedright (\textquotedbl{}language model\textquotedbl{} OR \textquotedbl{}language models\textquotedbl{} OR LLM OR LLMs OR \textquotedbl{}GPT-4\textquotedbl{} OR \textquotedbl{}zero-shot\textquotedbl{}) AND (screening OR \textquotedbl{}study selection\textquotedbl{} OR \textquotedbl{}citation screening\textquotedbl{} OR \textquotedbl{}abstract screening\textquotedbl{}) AND (\textquotedbl{}systematic review\textquotedbl{} OR \textquotedbl{}systematic reviews\textquotedbl{} OR \textquotedbl{}literature review\textquotedbl{}) AND (sensitivity OR recall OR specificity OR missed)\par}
\item[q4x, counter query] {\footnotesize\ttfamily\raggedright (LLM OR GPT OR ChatGPT) AND (screening) AND (unreliable OR inconsistent OR poor OR \textquotedbl{}not suitable\textquotedbl{} OR limitations OR \textquotedbl{}low sensitivity\textquotedbl{}) AND (\textquotedbl{}systematic review\textquotedbl{})\par}
\item[q4y, counter query] {\footnotesize\ttfamily\raggedright (\textquotedbl{}large language model\textquotedbl{} OR \textquotedbl{}large language models\textquotedbl{} OR \textquotedbl{}language model\textquotedbl{} OR LLM OR GPT OR ChatGPT) AND (\textquotedbl{}title and abstract\textquotedbl{} OR \textquotedbl{}abstract screening\textquotedbl{} OR \textquotedbl{}citation screening\textquotedbl{} OR \textquotedbl{}title screening\textquotedbl{} OR screening) AND (\textquotedbl{}systematic review\textquotedbl{} OR \textquotedbl{}systematic reviews\textquotedbl{}) AND (\textquotedbl{}low sensitivity\textquotedbl{} OR missed OR \textquotedbl{}false negatives\textquotedbl{} OR unreliable OR inconsistent OR \textquotedbl{}not replace\textquotedbl{} OR \textquotedbl{}cannot replace\textquotedbl{} OR poor)\par}
\item[q4z, counter query] {\footnotesize\ttfamily\raggedright (\textquotedbl{}title only\textquotedbl{} OR \textquotedbl{}titles only\textquotedbl{} OR \textquotedbl{}title-only\textquotedbl{} OR \textquotedbl{}title screening\textquotedbl{}) AND (screening) AND (missed OR \textquotedbl{}lower sensitivity\textquotedbl{} OR \textquotedbl{}false negatives\textquotedbl{} OR insufficient OR \textquotedbl{}relevant studies\textquotedbl{}) AND (\textquotedbl{}systematic review\textquotedbl{} OR \textquotedbl{}systematic reviews\textquotedbl{})\par}
\end{description}


\begin{table}
  \begin{sidecaption}{%
      Hits per query and database for the fourth question (screening
      with a language model), searched 2026-11-23 with Scholar, as
      \texttt{scholar sessions show} reports them; the databases and
      their syntax as in \cref{tab:stopping-matrix}.
      \enquote{of}: the database reported more records than it returned
      (\enquote{of about}: Semantic Scholar's approximate total).
    }[tab:llm-matrix]
    \footnotesize
    \begin{tabular}{@{}l*{6}{>{\raggedleft\arraybackslash}p{1.45cm}}@{}}
\toprule
Query & OA & S2 & Scopus & WoS & IEEE & arXiv \\
\midrule
q4a & 262 & 218 of about 220 & 197 & 94 & 13 & 24 \\
q4b & 247 & 202 of about 232 & 164 & 103 & 6 & 0 \\
q4c & 169 & 77 of about 78 & 88 & 43 & 9 & 16 \\
q4d & 801 & 397 of about 399 & 662 & 251 of 252 & 69 & 150 \\
q4e & 581 of 587 & 304 of about 306 & 210 & 174 & 21 & 47 \\
q4x & 196 & 136 & 116 & 85 & 14 & 19 \\
q4y & 195 & 123 & 84 & 89 & 9 & 9 \\
q4z & 51 of 52 & 44 of about 72 & 28 & 20 & 1 & 0 \\
\bottomrule
\end{tabular}

  \end{sidecaption}
\end{table}

Three studies were named before the search (an exploratory study of
screening with language models, an evaluation of GPT-4 for screening and
data extraction, and a study of automated screening for clinical
reviews); all came back.
The check by meaning returned 4 works that no session held; one was a
comparison of screening by title with screening by title and abstract,
which one query's wording should have returned, and it is cited below
although no session of this question holds it.
Of the 2,574 distinct records, 652 were screened in the first tier and 12
more that five or more databases had returned; 107 bore on the question
(63 supporting the claim, 44 qualifying or contradicting it).
Of the 1,910 records left unscreened, none among a random sample of 60
bore on the question (95\,\% interval 0 to 6.0\,\%).
The search found studies of two-stage and combined set-ups (a small model
first, a larger one on the uncertain records) but read them only in their
abstracts and cites none, so nothing is said here about such set-ups.
The search used no classifier; the screening was an agent's
(\cref{sec:search-screening}).
The records judged bearing are listed in \cref{sec:llm-hits}.

\section{Results}
\label{sec:llm-results}

\paragraph{High sensitivity on average, in health.}
A meta-analysis of 18 studies of screening medical literature with
language models, read as a preprint, reports that
\textcquote[preprint, abstract]{Xie2026}{\textins*{i}n title and abstract
screening, the pooled sensitivity was 0.92 and pooled specificity was
0.94}.
Over five reviews and 22,666 citations, for GPT-3.5 used to rule records
out, \textcquote[preprint, abstract]{Tran2024}{\textins*{p}ooled estimates
of sensitivity and specificity were 97.1\,\% (95\,\%CI 89.6\,\% to
99.2\,\%) and 37.7\,\%, (95\,\%CI 18.4\,\% to 61.9\,\%), respectively},
and limiting human screening to the records not ruled out could cut the
workload by 65\,\%.
A systematic review of such studies, read only in its abstract, found
that \textcquote[abstract]{Sandner2025}{\textins*{w}hile several
approaches achieve sensitivity above 95\,\%, none consistently reach the
99\,\% threshold required for replacing human screening}.
Small models, in a study read only in its abstract, were sensitive and
imprecise: \textcquote[abstract]{Sciurti2025}{\textins*{a}cross the
systematic reviews, LLMs achieved high sensitivity (up to 100\,\%) and
low precision (below 10\,\%) for records included by full text}, at
\$0.14 to \$1.93 per review for the model used through an interface, or
free for the two run locally.

\paragraph{Wide variation between reviews.}
On 500 abstracts of one systematic review with 12 included studies, the
best models were as accurate as human readers and still missed most of
the included studies:
\textcquote[preprint, abstract]{Krag2024}{GPT-4o, GPT-4T, and
Claude3-Opus achieved the highest accuracies (97\,\% to 98\,\%) comparable
to the human readers (96\,\% and 98\,\%), although sensitivity was low
(33\,\% to 50\,\%).}
Adjusted for chance agreement and the rarity of relevant records, GPT-4's
screening fared worse than its raw accuracy suggests:
\textcquote[abstract]{Khraisha2024}{barring studies that used highly
reliable prompts -- screening performance levelled at none to moderate
for different stages and languages}.
In software engineering, eight newer models did only slightly better than
seven older ones on a screening benchmark (mean Matthews correlation
0.365 against 0.347), differences between reviews exceeded those between
models, and the authors conclude, in a preprint, that
\textcquote[abstract]{Huotala2026}{LLMs are not yet ready to replace
humans in paper screening and the advantages new, more costly models
bring, appear to be very limited}.

\paragraph{People miss studies too, and runs differ.}
In a preregistered comparison within one complex scoping review, reported
in a preprint, two human workflows and two runs of one model reached a
similar recall of 82.3 to 82.9\,\%, none found every eligible record, and
\textcquote[abstract]{Figalova2026}{\textins*{t}wo nominally identical
GPT5.4 file-batch runs agreed on 91.7\,\% of records, yet differed on 94
records, including 29 verified eligible records retained by only one
run}.
In our own test, two identical runs of the classifier on 50 papers
differed on one kept-or-discarded verdict (\cref{sec:measured-classifier}).

\paragraph{Titles alone.}
In one nutrition review, \textcquote[abstract]{Teo2023}{57.9\,\% of
studies were correctly identified at the first-level screening stage
using the Ti/O approach versus 94.7\,\% by the Ti + Ab approach}, title
only against title and abstract.
In another review, read only in its abstract, screening titles first and
screening titles with abstracts both found every included article: \textcquote[abstract]{Mateen2013}{\textins*{t}he final
systematic review included 13 articles, all of which were identified by
both screening strategies}.

\section{Discussion and limitations}
\label{sec:llm-discussion}

Most of the evidence is from health reviews, several sources are
preprints (the meta-analysis and the study of GPT-3.5 in versions that
precede their published ones, which were not obtained), and models
change faster than studies of them appear.
Average sensitivity hides the spread between reviews, which is what a
single review faces: models as accurate as human readers missed half or
more of the included studies of one review that had few of them.
The one study in computing is a preprint.
Three sources were read only in their abstracts.
This search measured no error rate of a model of its own, since no
classifier screened it.

\section{Conclusion}
\label{sec:llm-conclusion}

Not as a general rule: on health datasets language models were sensitive
on average, but no approach reliably reached the sensitivity that
replacing human screening would need, sensitivity varied widely between
reviews, two identical runs of one model differed on eligible records,
and screening by title alone lost studies in one of two reviews.
For the procedure this means that a model's verdicts, like an agent's,
are unreviewed until a person confirms them; that a second screen and its
agreement are measured where the result matters; and that the records a
model discards are sampled like any other unscreened rest.

\section{Future work}
\label{sec:llm-future}

This search did not read in full the systematic review of sensitivity
and workload, the study of compact models, or the studies of two-stage
set-ups, and it used preprints for two sources whose published versions
exist.
The field lacks studies of model screening outside health, beyond the one
software-engineering benchmark, and of the spread between runs and
between reviews rather than averages; a study that measured both for
reviews in computing would say how large a second screen needs to be.

\section{The records judged bearing}
\label{sec:llm-hits}

% Author's note: table by generate_audit_table on a copy of screenlit-q4,
% bearing rows only, cited = cited in doc/bibliography.bib.
% --- scholar LaTeX fragment --------------------------------------------
% Generated by: scholar sessions export -f table
% Session: screenlit-q4
% \input this file from the body of a document whose
% preamble loads:
%   \usepackage{longtable}
%   \usepackage{booktabs}
%   \usepackage{array}
% ----------------------------------------------------------------------

\begingroup\footnotesize
\begin{longtable}{@{}%
  >{\raggedright\arraybackslash}p{0.44\textwidth}c%
  >{\raggedright\arraybackslash}p{0.13\textwidth}%
  >{\raggedright\arraybackslash}p{0.26\textwidth}@{}}
\caption{Records bearing on the claim, Question 4, screening with a language model (10 cited, 54 supporting and 34 qualifying or contradicting it, out of 2,478 unique hits; the 95 adjacent hits, the 507 false hits and 1,778 unclassified are not listed). The last column gives the reason for each inclusion; for rows decided by an agent the agent's confidence is given. A row is cited when the text cites the work, whoever decided on its records. Databases searched: OpenAlex (OA; not open access), IEEE Xplore (IEEE), Web of Science (WoS), Scopus, Semantic Scholar (S2), arXiv. Searched 2026-11-23.}\label{tab:screenlit-q4}\\
\toprule
Title & Year & Database & Reason \\
\midrule
\endfirsthead
\multicolumn{4}{@{}l}{\emph{\tablename~\thetable{} (cont.)}}\\
\toprule
Title & Year & Database & Reason \\
\midrule
\endhead
\midrule \multicolumn{4}{r@{}}{\emph{continued on next page}}\\
\endfoot
\bottomrule
\endlastfoot
Assessing the coverage of PubMed, Embase, OpenAlex, and Semantic Scholar for automated single-database searches in living guideline evidence surveillance: a case study of the international polycystic ovary syndrome guidelines 2023 & 2025 & Scopus & \emph{cited}, pending, unclassified \\
Can large language models replace humans in the systematic review process? Evaluating GPT-4's efficacy in screening and extracting data from peer-reviewed and grey literature in multiple languages & 2023 & arXiv, Scopus, OA, S2 & \emph{cited}: qualifies or contradicts the claim (0.60) \\
Compact large language models for title and abstract screening in systematic reviews: An assessment of feasibility, accuracy, and workload reduction & 2025 & Scopus, OA, WoS, S2 & \emph{cited}: supports the claim (0.60) \\
Evaluating human and LLM screening workflows in a conceptually complex scoping review: Recall--workload trade-offs and run-to-run consistency & 2026 & OA, arXiv, S2 & \emph{cited}: qualifies or contradicts the claim (0.60) \\
Evaluating Large Language Models for Literature Screening: A Systematic Review of Sensitivity and Workload Reduction & 2025 & OA, S2 & \emph{cited}: qualifies or contradicts the claim (0.70) \\
Has LLM Screening Performance Stalled in Software Engineering Systematic Reviews? & 2026 & arXiv, OA, S2 & \emph{cited}: qualifies or contradicts the claim (0.60) \\
Large language models for abstract screening in systematic- and scoping reviews: A diagnostic test accuracy study & 2024 & OA, S2 & \emph{cited}: qualifies or contradicts the claim (0.65) \\
Performance of Large Language Models in Automated Medical Literature Screening: A Systematic Review and Meta-analysis & 2026 & OA, WoS, Scopus & \emph{cited}: supports the claim (0.75) \\
Sensitivity and Specificity of Using GPT-3.5 Turbo Models for Title and Abstract Screening in Systematic Reviews and Meta-analyses & 2024 & Scopus, OA, WoS, S2 & \emph{cited}: supports the claim (0.70) \\
Title-plus-abstract versus title-only first-level screening approach: a case study using a systematic review of dietary patterns and sarcopenia risk to compare screening performance & 2023 & OA, WoS, Scopus, S2 & \emph{cited}: qualifies or contradicts the claim (0.70) \\
A comprehensive systematic review dataset is a rich resource for training and evaluation of AI systems for title and abstract screening & 2025 & OA, WoS, Scopus, S2 & supports the claim (0.50) \\
A Hybrid Semi-Automated Workflow for Systematic and Literature Review Processes with Large Language Model Analysis & 2024 & OA, WoS, Scopus, S2 & supports the claim (0.50) \\
A prompt-based tutorial for large language model-assisted screening in systematic reviews and meta-analyses & 2026 & WoS, Scopus, OA, S2 & supports the claim (0.50) \\
A Reproducible LLM-Ensemble Protocol for Title--Abstract Screening in OR/MS Systematic Reviews: Design and Proof of Concept & 2026 & OA, S2 & supports the claim (0.50) \\
Artificial intelligence-driven study selection in systematic reviews of randomized controlled trials, emulated trials and economic evaluation studies using large language models & 2026 & Scopus, OA, WoS, S2 & supports the claim (0.50) \\
Assessing the Ability of ChatGPT to Screen Articles for Systematic Reviews & 2023 & arXiv & supports the claim (0.50) \\
Automated full-text screening and accelerated reviews using large language models with context-aware agents: an exploratory analysis in biomarker research & 2026 & OA, WoS, Scopus, S2 & supports the claim (0.50) \\
Automated Paper Screening for Clinical Reviews Using Large Language Models: Data Analysis Study & 2024 & OA, WoS, Scopus, arXiv, S2 & supports the claim (0.50) \\
Automating Screening of Titles and Abstracts in Systematic Reviews: An Assessment of GPT-4o mini & 2026 & OA, S2 & supports the claim (0.50) \\
Collaborating with large language models in literature screening for a systematic review of college students' GenAI literacy & 2026 & OA, WoS, Scopus, S2 & supports the claim (0.50) \\
Collaborative large language models (LLMs) are all you need for screening in systematic reviews & 2026 & OA, S2 & supports the claim (0.50) \\
Comparative Assessment of ChatGPT, DeepSeek, and Human Reviewers for Full-Text Screening in Systematic Reviews on the Impact of Air Pollution in Respiratory Diseases & 2025 & OA, WoS, Scopus & supports the claim (0.50) \\
Comparing Open-Weight Large Language Models for Systematic Review Tasks: One Health Antimicrobial Resistance Case Study (Preprint) & 2026 & OA & supports the claim (0.50) \\
Cost-Efficient Large Language Model-Assisted Title and Abstract Screening for Systematic Reviews: Method Evaluation and Validation & 2026 & OA & supports the claim (0.50) \\
Development and evaluation of prompts for a large language model to screen titles and abstracts in a living systematic review & 2025 & OA, WoS, Scopus, S2 & supports the claim (0.50) \\
Development of meta-prompts for Large Language Models to screen titles and abstracts for diagnostic test accuracy reviews & 2023 & OA, S2 & supports the claim (0.50) \\
Development of Prompt Templates for Large Language Model--Driven Screening in Systematic Reviews & 2025 & Scopus, OA, WoS, S2 & supports the claim (0.50) \\
Dual-Model LLM Ensemble via Web Chat Interfaces Reaches Near-Perfect Sensitivity for Systematic-Review Screening: A Multi-Domain Validation with Equivalence to API Access & 2025 & OA, S2 & supports the claim (0.50) \\
Efficiency of LLM-Based Systematic Review Screening: Mimicking Human Expert Cognition & 2025 & OA & supports the claim (0.50) \\
Ensembles of Large Language Models for Identifying EQ-5D Studies in PubMed Based on Their Abstracts & 2026 & arXiv, OA, IEEE, Scopus, S2 & supports the claim (0.50) \\
ENTGPT: Applying Large Language Models to Systematic Review Screening With the Novel STARR Protocol & 2026 & Scopus, OA, WoS, S2 & supports the claim (0.50) \\
Evaluating a Locally Deployed 20-Billion Parameter Large Language Model for Automated Abstract Screening in Systematic Reviews & 2027 & WoS, Scopus, OA, S2 & supports the claim (0.50) \\
Evaluating large language models for title/abstract screening: a systematic review and meta-analysis \& development of new tool & 2025 & OA, WoS, Scopus, S2 & supports the claim (0.65) \\
Evaluating the Efficacy and Efficiency of GPT-5 for Automated Title and Abstract Screening in Orthopedic Surgery Systematic Reviews & 2025 & OA, WoS, Scopus, S2 & supports the claim (0.50) \\
Evaluating the Efficacy of Large Language Models for Systematic Review and Meta-Analysis Screening & 2024 & OA, S2 & supports the claim (0.50) \\
Expediting citation screening using PICo-based title-only screening for identifying studies in scoping searches and rapid reviews & 2017 & OA, WoS, Scopus, S2 & supports the claim (0.45) \\
Fine-Tuning A Large Language Model for Systematic Review Screening & 2026 & OA, arXiv, S2 & supports the claim (0.50) \\
GPT-3.5 Turbo and GPT-4 Turbo in Title and Abstract Screening for Systematic Reviews & 2025 & WoS, Scopus, OA, S2 & supports the claim (0.50) \\
High-performance automated abstract screening with large language model ensembles & 2024 & arXiv, OA, WoS, Scopus, S2 & supports the claim (0.50) \\
How Well Do ChatGPT and Claude Perform in Study Selection for Systematic Review in Obstetrics & 2025 & OA, WoS, Scopus, S2 & supports the claim (0.50) \\
Human-Comparable Sensitivity of Large Language Models in Identifying Eligible Studies Through Title and Abstract Screening: 3-Layer Strategy Using GPT-3.5 and GPT-4 for Systematic Reviews & 2024 & OA, WoS, Scopus, S2 & supports the claim (0.50) \\
Large language model-based multiagent collaboration for abstract screening toward automated systematic reviews & 2026 & WoS, Scopus, OA, S2 & supports the claim (0.50) \\
Leveraging LLMs for Title and Abstract Screening for Systematic Review: A Cost-Effective Dynamic Few-Shot Learning Approach & 2025 & arXiv, OA, WoS, Scopus, S2 & supports the claim (0.50) \\
Lightweight open-source reasoning large language model for evidence synthesis screening & 2026 & OA, WoS, Scopus, S2 & supports the claim (0.55) \\
Making Broad Evidence Synthesis Feasible: An LLM Screening Agent for Meta-Analyses Applied To Suicide Prevention & 2026 & OA, S2 & supports the claim (0.50) \\
MechaScreener: Large Language Model-Based Automated Screening for Systematic Reviews and Research & 2026 & OA, Scopus, S2 & supports the claim (0.50) \\
MechaScreener: Large Language Model-Based Automated Screening for Systematic Reviews and Research & 2026 & OA & supports the claim (0.50) \\
Methodological insights into ChatGPT's screening performance in systematic reviews & 2024 & OA, WoS, Scopus, S2 & supports the claim (0.50) \\
Operationalising LLM-assisted screening of literature to support systematic reviews & 2026 & OA, S2 & supports the claim (0.50) \\
Optimal large language models to screen citations for systematic reviews & 2025 & OA, WoS, Scopus, S2 & supports the claim (0.50) \\
Performance Comparison of Large Language Models for Efficient Literature Screening & 2025 & OA, WoS, Scopus, S2 & supports the claim (0.50) \\
Performance of a Large Language Model in Screening Citations & 2024 & Scopus, OA, WoS, S2 & supports the claim (0.50) \\
Qualitative Analysis of Discrepancy Patterns Between Large Language Models and Human Reviewers in Abstract Screening for Systematic Reviews & 2026 & OA, WoS, S2 & supports the claim (0.50) \\
REFLECTIVE-TIAB: cost-effective prompt optimization for large language model-based title and abstract screening in literature reviews & 2026 & OA, WoS, Scopus, S2 & supports the claim (0.50) \\
Synapse: Automated Generation of Systematic Review Protocols and AI-Assisted Screening with Large Language Models & 2026 & OA, S2 & supports the claim (0.50) \\
Towards Automating Articles Screening Processes Using Chain-of-Thought Large Language Models & 2026 & OA, WoS, Scopus, S2 & supports the claim (0.50) \\
Transforming literature screening: The emerging role of large language models in systematic reviews & 2025 & Scopus, S2 & supports the claim (0.50) \\
Usefulness of machine learning softwares to screen titles of systematic reviews: a methodological study & 2023 & OA, WoS, Scopus, S2 & supports the claim (0.50) \\
Using a large language model as a third reviewer to augment dual human fulltext screening in orthopaedic systematic reviews & 2026 & OA, WoS, Scopus, S2 & supports the claim (0.50) \\
Using full agreement across multiple large language models for title-and-abstract screening in systematic reviews: a proof-of-concept & 2026 & Scopus, S2 & supports the claim (0.50) \\
Using large language models to directly screen electronic databases as an alternative to traditional search strategies such as the Cochrane highly sensitive search for filtering randomized controlled trials in systematic reviews & 2025 & Scopus, WoS, OA, S2 & supports the claim (0.50) \\
Validation of Generative AI Models to Expedite Title and Abstract Screening in Systematic Reviews & 2025 & OA & supports the claim (0.50) \\
Validation of large language models (Llama 3 and ChatGPT-4o mini) for title and abstract screening in biomedical systematic reviews & 2025 & Scopus, OA, WoS, S2 & supports the claim (0.50) \\
Zero-shot Generative Large Language Models for Systematic Review Screening Automation & 2024 & arXiv, OA, WoS, Scopus, S2 & supports the claim (0.50) \\
A biomedical BERT ensemble (TITAN-SR) outperformed active-learning and large language model chatbots for systematic review title and abstract screening: development and external validation & 2026 & OA, WoS, Scopus, S2 & qualifies or contradicts the claim (0.50) \\
A Machine Learning and Large Language Model Tool for Systematic Literature Reviews of Health Economic Evidence: A Validation Study & 2026 & OA, WoS, Scopus, S2 & qualifies or contradicts the claim (0.50) \\
Automated title and abstract screening for scoping reviews using the GPT-4 Large Language Model & 2023 & OA, arXiv, S2 & qualifies or contradicts the claim (0.50) \\
Batch Size Effects on Mid2025 StateoftheArt Large Language Model Performance in Automated Title and Abstract Screening & 2026 & OA & qualifies or contradicts the claim (0.50) \\
Beyond Accuracy: LLM Variability in Evidence Screening for Software Engineering SLRs & 2026 & arXiv, OA, S2 & qualifies or contradicts the claim (0.55) \\
Cal-X: Enhancing Systematic Review Screening with LLMs and Next-Token Likelihood Calibration & 2025 & IEEE, Scopus, OA, S2 & qualifies or contradicts the claim (0.50) \\
ChatGPT (GPT-3.5) versus human reviewers for title-and-abstract screening in a scoping review: A retrospective comparative evaluation & 2026 & WoS, Scopus, OA, S2 & qualifies or contradicts the claim (0.50) \\
Class Imbalance and Batch Effects in LLM-Based Screening for Systematic Reviews & 2026 & OA, arXiv, S2 & qualifies or contradicts the claim (0.45) \\
Efficacy of Large Language Models for Screening of Systematic Reviews on Periprosthetic Joint Infection & 2026 & OA, WoS, Scopus, S2 & qualifies or contradicts the claim (0.50) \\
Evaluating and explaining LLM classification for automatic screening in Systematic Reviews & 2026 & Scopus, OA, S2 & qualifies or contradicts the claim (0.50) \\
Evaluating Large Language Models for Radiology Systematic Review Title and Abstract Screening & 2025 & OA, WoS, Scopus, S2 & qualifies or contradicts the claim (0.50) \\
Generative artificial intelligence use in evidence synthesis: A systematic review & 2025 & OA, WoS, Scopus, S2 & qualifies or contradicts the claim (0.60) \\
How Far Can a Language Model Carry a Systematic Review? Preregistered Reliability Evidence From Screening to Measurement-Quality Coding & 2026 & OA & qualifies or contradicts the claim (0.55) \\
Is Large Language Model-Assisted Citation Screening Feasible in a Scoping Review on Nonpharmacological Interventions for Delirium in Patients With Cancer? & 2025 & OA, S2 & qualifies or contradicts the claim (0.50) \\
Large Language Model-Assisted Title and Abstract Screening for A Scoping Review of Circulating miRNA Monitoring: Pilot Validation Against a Human Reviewer and Quantifying the Error of Keyword Pre-Filters & 2026 & OA & qualifies or contradicts the claim (0.50) \\
Large language models for conducting systematic reviews: on the rise, but not yet ready for use---a scoping review & 2025 & OA & qualifies or contradicts the claim (0.55) \\
Large language models for literature screening in conceptually complex and interdisciplinary reviews & 2026 & OA, WoS, Scopus, S2 & qualifies or contradicts the claim (0.50) \\
Large language models for the screening step in systematic reviews in dentistry & 2025 & OA, WoS, Scopus, S2 & qualifies or contradicts the claim (0.50) \\
Large Language Models for Title/Abstract Screening in Systematic Literature Reviews: A Case Study in Precision Livestock Farming & 2025 & OA, IEEE, Scopus, S2 & qualifies or contradicts the claim (0.50) \\
Large language models show promising performance for some systematic review tasks but call for cautious implementation: a systematic review & 2026 & OA, WoS, Scopus, S2 & qualifies or contradicts the claim (0.50) \\
Linguistic Polarity and Decision Architecture in LLM-Based Abstract Screening for Systematic Reviews & 2026 & Scopus, OA, WoS, S2 & qualifies or contradicts the claim (0.50) \\
LLM-assisted screening and extraction for systematic reviews: workload reduction and refusal-driven missingness in a land-use evidence synthesis & 2026 & OA & qualifies or contradicts the claim (0.55) \\
Performance of large language models conducting systematic review tasks in prosthodontics & 2026 & OA, WoS, Scopus, S2 & qualifies or contradicts the claim (0.50) \\
Performance of LLMs in Citation Screening: A Comparison Across Datasets with Varied Inclusion Rates & 2025 & Scopus, OA, WoS, S2 & qualifies or contradicts the claim (0.55) \\
SESR-Eval: Dataset for Evaluating LLMs in the Title-Abstract Screening of Systematic Reviews & 2025 & arXiv, OA, IEEE, WoS, Scopus, S2 & qualifies or contradicts the claim (0.50) \\
The AI Reviewer: Evaluating AI's Role in Citation Screening for Streamlined Systematic Reviews & 2025 & OA, WoS, Scopus & qualifies or contradicts the claim (0.50) \\
The Promise and Challenges of Using LLMs to Accelerate the Screening Process of Systematic Reviews & 2024 & OA, arXiv, WoS, Scopus, S2 & qualifies or contradicts the claim (0.50) \\
TiAb Review Plugin: A Browser-Based Tool for AI-Assisted Study Selection in Systematic Reviews & 2026 & arXiv, OA, S2 & qualifies or contradicts the claim (0.50) \\
Title and abstract screening for literature reviews using large language models: an exploratory study in the biomedical domain & 2024 & OA, WoS, Scopus, S2 & qualifies or contradicts the claim (0.60) \\
Understanding LLMs in Title-Abstract Screening: From Disagreements to Recommendations & 2027 & Scopus, arXiv, S2 & qualifies or contradicts the claim (0.50) \\
Using GPT4 for Title and Abstract Screening in a Literature Review of Public Policies: A Feasibility Study & 2025 & OA, S2 & qualifies or contradicts the claim (0.50) \\
Validating Large Language Models for Title-Abstract Screening in Low-Prevalence Systematic Reviews: An Environmental Science Case Study & 2026 & OA, WoS, Scopus, S2 & qualifies or contradicts the claim (0.55) \\
Validation of a two-stage automated screening pipeline for medical systematic reviews: a stratified agreement study with three independent human reviewers, using PABAK and Gwet's AC1 & 2026 & OA, S2 & qualifies or contradicts the claim (0.50) \\
Validation of an LLM-assisted title and abstract screener on two published systematic reviews & 2026 & OA & qualifies or contradicts the claim (0.60) \\
\end{longtable}
\endgroup

